The \emph{geometrical} Volume-of-Fluid (VoF) method is a method where the 
interface is represented by a set of simple geometrical entities and advection 
of the volume fraction is accomplished using geometrical algorithms. When
the interface is represented as a set of polygons (planes) reconstructed  from
the volume fraction field, the geometrical VoF algorithm is commonly known as
the Piecewise Linear Interface Calculation (PLIC) algorithm \citep{deBarPLIC,
RiderReconstructing1998}. Geometrical operations involving interface planes and
mesh cells are then performed in order to advect the volume fraction field.

The widespread use of the geometrical Volume of Fluid (VoF) method for
simulating two-phase flow is mostly based on its intrinsic mass conservation 
and the ability to automatically deal with topological changes of the fluid 
interface (no special operations on the discrete interface geometry are required). 
However, the majority of the developments related to the geometrical VoF method 
have so far been done on Cartesian and structured meshes because of the high 
numerical accuracy they provide and the simplicity of the algorithms relying 
on structured mesh points. Problems involving complex geometrical shapes of the 
flow domain, usually rely on automatic (fast) algorithms for generation of 
unstructured meshes. For this reason, the use of unstructured meshes has prevailed 
in those parts of the industry and research where complex geometry of the flow 
domain is of importance. However, the use of an accurate geometrical VoF method 
for such problems was not possible so far; instead VoF methods involving algebraic 
advection of the volume fraction field are in widespread use. 

This contribution utilises the OpenFOAM{\textsuperscript{\textregistered}}
software for Computational Continuum Mechanics (CCM) as an implementation
platform for the geometrical VoF method on arbitrary unstructured meshes. As 
detailed in the following, the developed geometrical VoF method is based on novel 
a directionally unsplit advection algorithm using a full discrete Lagrangian flow 
map on unstructured meshes and thus allowing for accurate direct numerical 
simulation (DNS) of two-phase flow in domains of arbitrary geometrical complexity. 
The geometrical VoF algorithm for volume fraction transport is shown to be robust, 
volume conservative and boundedness preserving; it is fully parallelized using 
the domain decomposition approach and supports dynamic local Adaptive Mesh 
Refinement (AMR). The algorithm utilizes a new efficient (quadratic convergent) and 
accurate iterative flux correction algorithm, that avoids topological changes of 
the flux polyhedra and inherently ensures mass conservation. To the best of the 
authors' knowledge, for the first time such a geometrical VoF method is shown to 
be capable to deal with three dimensional, spatially complex and time varying 
velocity fields. The obtained volume of symmetric difference interface 
reconstruction errors are the lowest reported so far in the literature for 
unstructured meshes. %
The development of the geometrical VoF algorithm significantly extends the method 
base of the OpenFOAM library using a modern policy based design for high program 
code modularity and combinatorial complexity of sub-algorithms for geometrical 
reconstruction and advection. %

OpenFOAM is a Finite Volume Method (FVM) library for Computational Continuum 
Mechanics (CCM) written in the C++ programming language. It has a modular structure 
and provides a possibility of an extensive re-use of the available functionality. 
A broad choice of mathematical and constitutive models is available, as well as 
linear solvers, dynamic mesh operations involving both mesh motion and topological
operations and a comprehensive set of boundary conditions, etc. Details of the 
development model of OpenFOAM and its relationship to continuum mechanical modeling 
and simulations of CCM problems are provided by \citet{WellerJasak1998} and
\citet{Jasak2007}. The benefit of developing the geometrical VoF method for 
unstructured  meshes using OpenFOAM within an \emph{Open Source} framework is thus 
two-fold: it provides an implementation with the ability to deal with arbitrarily 
complex flow domains, and it allows for extensions of the method by coupling with 
the present functionality of the OpenFOAM library. 

The goal of this work is to develop a geometrical VoF library and a set of
solver applications that can deal with two-phase flows involving complex
transport processes (\emph{complex flows}) in complex domains (\emph{complex
geometry}). The topology of an unstructured meshes dictates our choice of 
algorithms for both the reconstruction of the interface geometry and the advection 
of the volume fraction, as well as the underlying geometrical operations. The fluid
interface is approximated with a piecewise planar geometry which is
reconstructed using the method of \citet{ParkerYoungsGradient} with an
alternative accurate parallelized gradient calculation on arbitrary
unstructured mesh and is able to deal with arbitrary shapes of the mesh cells.
An Eulerian unsplit geometrical algorithm is chosen for the volume fraction
advection, based on a full unique Lagrangian trajectories located in the cell
corner points. We do not consider complicated geometrical operations necessary
for non-convex congruent polyhedra, which may be constructed by the advection
algorithm in those parts of the flow domain where insufficient resolution is
provided. The reason for this decision is that the flow will
be sufficiently resolved by local dynamic AMR near the interface. 
 
The accuracy of the geometrical VoF method can be enhanced
when coupled with dynamic local Adaptive Mesh Refinement (AMR)\footnote{From
this point on in the text, dynamic local AMR will be refered to as AMR since it
must follow an evolving fluid interface when applied to two-phase flow
simulations.} (see \citet{KotheUnstructured1996} and \citet{Cerne2002} for
details). Extensions to higher order geometrical methods such as the Moment-of-Fluid 
(MoF) method of \citet{AhnShashkov2009} or higher order reconstruction
(e.g.  Patterned Interface Reconstruction by \citet{MossoPIR2009}) are also
possible if it proven to be necessary (even with applied AMR). 


\subsection{State of the art}

For the sake of clarity, we provide a brief overview of the important aspects
of the geometrical VoF method as well as the related terminology. Emphasis is
put upon geometrical VoF methods and related algorithms rather than the algebraic
VoF method. The interested reader is referred to appropriate references in 
literature once the scope of this overview is left. 

In order to advect the volume fraction field using geometrical algorithms, the
geometrical representation of the interface needs to be sufficiently accurate,
and the geometrical advection algorithm must conserve mass as well as avoid
artificial deformation of the interface, while keeping the volume fraction
field bounded between values of $0$ and $1$. 

Since the interface is described as a set of planes (polygons) the orientation 
of which is given by the volume fraction field, the planes are disconnected
on cell faces. Discontinuities of the interface planes at cell faces may be 
described in terms of \emph{geometrical stability} of the interface.
The discontinuities decrease with the increase of mesh
resolution and the accuracy of the gradient calculation, used to determine the
orientation of the plane normals. Furthermore, the orientations of the
interface planes may be additionally optimized in a way that does not impose
errors in mass conservation. Such second order convergent reconstruction
algorithms\footnote{Order of convergence of geometrical algorithms is related
to convergence of standard geometrical reconstruction or advection errors.}
result in smaller discontinuities, but involve additional computational
complexity. However, it is of utmost importance to reduce them to a minimum, 
since geometrical instabilities generated by the reconstruction algorithm cause
accumulation of other errors during the advection of the volume fraction field
\citep{Cerne2002}. 

Geometrical advection algorithms may be sorted into two general categories
based on the direction splitting: directional (operator) split and un-split
algorithms. Another division may be defined based on the choice of the mesh
geometry used to advect the volume fraction field: \emph{face swept}, and
\emph{Lagrangian re-mapping} algorithms. 

\emph{Operator split algorithms} compute the volume fraction fluxes
interchangeably in directions orthogonal to the coordinate axes. Such ordering
is native to a structured mesh, but can be extracted from the unstructured
hexahedral mesh as well. Since the advection directions are interchanged, each
advection iteration is followed by a reconstruction step where the interface
geometry is computed from the updated volume fraction field. 

\emph{Directional un-split algorithms} compute the advection of the volume
fraction field in a single computational step, using the velocity vectors.
Hence, a single reconstruction of the interface is required after each
advection step.

\emph{Face-swept algorithms} rely on flux calculations coming from sweeping or
\emph{back-tracing} cell faces along the streamlines given by the velocity
field. A \emph{swept polyhedron} is then adjusted, such that its volume equals
the value of the volumetric flux coming from the numerical solution of the
pressure equation in order to assure mass conservation. 

\emph{Eulerian backtracking / Lagrangian re-mapping} algorithms compute the
increase of the volume fraction on a per-cell basis, by sweeping the mesh
backward using the velocity field interpolated to the mesh vertices, and
intersecting the swept mesh with the geometrical data of the interface. After
this Lagrangian back-tracking step, a re-mapping is executed on the underlying
Eulerian mesh to set the new volume fraction values.

On unstructured meshes, both face-swept and Lagrangian-Eulerian re-mapping
algorithms are directionally un-split. Depending on the complexity of the flow
field, both face-swept and Lagrangian re-mapping algorithms encounter problems in
geometrical operations because the created polyhedra have non-planar faces and
may even be non-convex.  

A detailed review of algorithms involving two-phase flow methods on structured
meshes with emphasis on the geometrical VoF method can be found in a book on
Direct Numerical Simulations (DNS) of two-phase flows by
\citet*{Tryggvason2011direct} and in journal contributions by
\citet{RiderReconstructing1998}, \citet{ScardovelliReview1999},
\citet{ScardovelliZaleskiEILE2003}, and \citet{Aulisa2007}. For latest 
developments on Cartesian structured meshes the interested reader is 
referred to \citep{WeymouthYue2010,MencingerZunAdaptive2011,LeChenadec2013} and 
the references therein. %
A detailed numerical analysis of the errors native to the two-dimensional 
geometrical VoF method on structured meshes is described in \citet{Cerne2002}. 
So far, main contributions to geometrical algorithms for VoF methods were 
concentrated on the following aspects: improving interface normal orientation, 
improving overall mass conservation, reducing under-and-overshoots of the volume 
fraction, and removing wisps and artificial interface separation. 

In the remainder we present a brief overview of the research efforts involving 
three-dimensional calculations and geometrical VoF algorithms which support 
\emph{unstructured} domain discretization, since they are more relevant to our work. 

\citet{Swarz1996} and \citet{KotheUnstructured1996} have produced first (to our
knowledge) publications referring to a geometrical VoF method with support for
unstructured meshes. 

\citet{KotheUnstructured1996} emphasize that the motivation behind research
towards unstructured meshes is based on the requirement for arbitrary
geometrical complexity of the flow domain. Important aspects of the gradient
reconstruction for the volume fraction field is presented, resulting in a
conclusion that a wide cell stencil ($27$ cells for a hexahedral mesh) is
necessary in order to obtain geometrical stability of the reconstructed
interface. A valuable note is provided regarding the visualization of results,
which emphasizes the importance of viewing the actual polygonal geometry of the
reconstructed interface. Standard visualization practice utilises an
iso-surface representation of the interface, where the iso-value of the volume
fraction is prescribed (e.g. to a value of $0.5$). However, the use of
iso-contours for visualizing interfaces introduces artificial connectivity of
the interface e.g. for thin filaments. Besides this, the iso-contour hides all
the errors which are shown when the polygonal geometry is used to visualize the
interface: \emph{wisps} -- small differences in the volume fraction in
the bulk of the phases from values $0$ or $1$ that cause reconstructed
interface geometry to appear (cp.\ \citet{HarvieFletcherStream2000} for
details) -- or even larger artificially separated parts of the interface 
may not be shown.

\citet{Swarz1996} developed a second order convergent interface reconstruction
method for unstructured meshes. The reconstruction is based on optimizing the
orientation of the plane normals taking into account orientation of the planes
in the neighbouring interface cells. For the advection of the volume fraction,
a full Lagrangian back tracing of the finite volume mesh is performed, followed
by a geometrical intersection with the PLIC geometry, and remapping of the
volume fraction field onto the Eulerian background mesh. This advection
algorithm is directionally un-split, it relies completely on the discrete flow
map and requires a single reconstruction step per advection iteration. 

\citet{KotheVoFchar1999} provides an overview of the geometrical VoF method
developed for simulating metal casting processes in complex molds. Although a
casting process is a specialized application of the geometrical VoF method, the
problems involving the complex domain geometry and physical characteristics of
the model (e.g. significant differences in values of physical properties
accross the interface), make the requirements applicable to a general purpose
geometrical VoF method. They show preliminary results involving reconstruction
of an interface on an unstructured tetrahedral mesh.

\citet{AhnShashkovGeometricAlg2008} present geometrical algorithms used for the
reconstruction and advection step operations, with support for polyhedral
cells. They report successful three-dimensional reconstructions of complex
geometrical shapes using the method of \citet{ParkerYoungsGradient} and the
Least-squares-Volume-of-fluid-Interface-Reconstruction-Algorithm (LVIRA)
developed by \citet{PuckettLvira1991}, both modified in order to support
polyhedral cell shapes. \citet{ShashkovDyadechkoMoF2005} present three
dimensional results of the Moment of Fluid (MoF) reconstruction method which
optimizes the difference in position of the barycentric centre of the truncated
polyhedron filled with a specific phase. The MoF method can be viewed as an
extension of the geometrical VoF method by introducing additional information
in the form of barycentric centres of the truncated cells, which results in a
more accurate reconstruction algorithm (see
\citet{AhnShashkovMultiMaterialRec2007} for details on multi-material
reconstruction on polyhedral meshes). 

\citet{jofreUnstructuredPLIC} show three-dimensional results for both
geometrical reconstruction and the volume fraction advection on unstructured
finite volume mesh. The authors present the volume fraction evolution for
rotational flow using a full discrete Lagrangian flow map. Obtained errors of
the advection algorithm are close in magnitude to those reported by other
authors on structured meshes (see \citet{LiovicCVTNA2006} and
\citet{HernandezLopezPart1Advection2008}). 

\citet{MossoPIR2009} developed the geometrical reconstruction algorithm called
Patterned Interface Reconstruction (PIR) algorithm. The algorithm is of second
order of convergence (reconctructs planes exactly) on tetrahedral (and
triangular) meshes. The interface geometry is initially reconstructed using the
Youngs PLIC reconstruction algoritm followed by linear or spherical smoothing
operations based on the geometrical pattern formed by a group of interface
polygons. Apart from reconstruction algorithm of the MoF method, this is the
only reconstruction algorithm of second order that directly applies to 
arbitrary unstructured meshes. 

\citet{LiovicCVTNA2006} have developed a three-dimensional extension of the
original algorithm of Rider and Kothe on Cartesian meshes with a new accurate
reconstruction algorithm and an elegant flux divergence correction method used
for improving the mass conservation. Errors in mass conservation appear because
of overlaps between the flux polyhedra on common mesh edges. Although the
algorithm shows good results, it relies on a directed connectivity between
cells (e.g. \emph{L-shaped stencils}) that requires the use of a structured
mesh. Additional computation of directed cell connectivity would be
possible for unstructured \emph{hexahedral} meshes, but not on arbtirary
unstructured meshes.

\citet{LopezEMFPA2004} show results with a full discrete flow map as a source
for the back tracing of the flux polyhedra on a two-dimensional Cartesian mesh.
Local geometrical corrections of the flux polyhedron are applied in order to
ensure mass conservation. For the reconstruction algorithm, a compact spline
reconstruction algorithm, with splines based on the centre points of the
reconstructed interface line segments is used. This work is extended to three
dimensions, again on a Cartesian mesh, by
\citet{HernandezLopezPart1Advection2008} and \citet{LopezPart2reconstruction}.
However, the construction of swept polyhedra is based on edge centres of the
Cartesian mesh, with velocities interpolated from the staggered velocity field
mapped to the face centres. Although a notable improvement is achieved, there
is still a possibility of an overlap between the flux polyhedra around face
corners. However, the mass conservation is strictly met by an analytical
approach to adjusting the swept polyhedron volume. 

\citet{IveyMoin2012} have implemented the so-called edge-matched flux polyhedra 
algorithm in 3D (EMFPA 3D), suggested as an improvement of the unsplit VoF 
algorithm developed by \citet{HernandezLopezPart1Advection2008}. %
The algorithm locally approximates a stream tube emanating from the flux face. 
The flux face is triangulated to construct the flux polyhedron, which then 
is modified to correct the non-planarity of the faces. Divergence correction, as 
recommended by \citet{HernandezLopezPart1Advection2008}, relies on an analytical 
expression for computing the volumes of convex polyhedrons. Adjusting the flux 
polyhedron introduces overlapping of the flux polyhedra that needs to be further 
corrected to prevent the volume fraction from over-/undershooting. Unfortunately, 
a validation of the results showing quantitative information on the standard 
advection errors are not provided.

%A conservative advection scheme based on the use of edge-matched flux polyhedra to integrate the volume fraction evolution equation on general grids is presented. The algorithm prevents the formation of over/undershoots of the volume fraction by enforcing that the flux polyhedra do not over/underlap, removing the need for unphysical and inaccurate redistribution algorithms.

%The method uses the flux polyhedron construct to locally approximate a stream tube emanating from the flux face, perturbing the definition in order to remove all over-/underlap to prevent the volume fraction from over-/undershooting.

Summarizing, all of the aforementioned algorithms rely on \emph{divergence corrections}  
in some form to account for the fact that the discrete velocity field used for the
construction of the flux polyhedron is not divergence free. Besides this issue,
wisps (see \citet{HarvieFletcherStream2000} for details) are present in the
volume fraction field, and are subsequently redistributed. 

As for AMR, \citet{PopinetTreeAMR2003} developed a very advanced geometrical VoF method 
that supports both 2D and 3D calculations with AMR on structured mesh. The
domain discretization is based on an \emph{octree} data structure which enables
fast and flexible mesh refinement operations and locally accurate
differentiation.  Recently, \citet{PopinetBalancedForce2009} has shown that the
AMR plays a significant role in reducing \emph{parasitic currents}. Details on
the capabilities and applications of the octree based geometrical VoF method
are provided by \citet{FusterPopinet2009} and \citet{AgbaglahPopinet2011}. 

Fully unstructured AMR available in OpenFOAM delegates the refinement to each
cell, and is hence able to deal in a straightforward way with arbitrary cell
shapes and solution domains of arbitrary geometrical complexity. With fully
unstructured AMR, a new mesh is created with refined and unrefined regions each
time a topological operation is executed. The octree based AMR implies a flow
domain in the shape of a regular hexahedron. %

We have chosen the fully unstructured AMR since it does not restrict us to a 
specific cell shape (allows automatic mesh generation of complex solution 
domains) and it is already a part of the OpenFOAM library, which means it can 
be coupled to other dynamic mesh functionality (mesh motion, mesh deformation, 
sliding interfaces, moving reference frames, etc.). 

In the following section the governing equations of the flow are described.
Geometrical operations required by the unstructured meshes and their use in the
reconstruction and advection step of the method are presented in the third
section. In the fourth section the coupling between the numerical solution of
the governing equations and the geometrical evolution of the volume fraction is
set out. Results and discussions are presented in the two final sections.
